\scanpage{189}
% manuscript position: top
\mspage{105}
% manuscript position: 1–3
d'où en appliquant $g$, qui commute aux $\varinjlim$ de type $I$, et se rappelant maintenant que $g\mid C\simeq f$. \deleted{et que $f$}
% manuscript position: 4
\[
g\big(\varinjlim_J{}_{\operatorname{Ind}_\pi C}\varphi(j)\big)
\simeq\varinjlim_{i\in I}f(\Psi(i)).\tag*{(\(\ast\ast\ast\))}
\]

% manuscript position: 5
On a d'autre part
% manuscript position: 6
\[
\Psi(i)\overset{\mathrm{déf}}{=}\varinjlim_{j\in J}\Phi(i,j)
\qquad\varinjlim\ \text{dans }C
\]
% manuscript position: 7–8
et comme $f$ commute aux $\varinjlim$ de type $J$ (car $\operatorname{card}\operatorname{Fl} J\leq\pi$), on en conclut
% manuscript position: 9
\[
f(\Psi(i))\simeq\varinjlim_{j\in J}f(\Phi(i,j)),
\]
% manuscript position: 10
donc par $(\ast\ast\ast)$
% manuscript position: 11
\[
g\big(\varinjlim_{j\in J}\varphi(j)\big)\simeq\varinjlim_{i\in I}\varinjlim_{j\in J}f(\Phi(i,j)).\tag*{(\(\ast\ast\ast\ast\))}
\]

% manuscript position: 12–14
En comparant avec $(\ast)$ et utilisant à nouveau \unsure{a|à} la commutation de $\varinjlim_I$ à $\varinjlim_J$, on trouve donc
% manuscript position: 15
\[
g\big(\varinjlim\varphi(j)\big)\xleftarrow{\approx}\varinjlim g(\varphi(j)),
\]
% manuscript position: 16–22
ce qui \deleted{\unsure{achève de}} prouve que $g$ \deleted{est ds \unsure{$\underline{\mathrm{Hom}}$}} commute aux limites inductives de type $J$, et achève de prouver que $\unsure{O|g|0}$ est $\underline{\mathrm{Hom}}_{!}(\operatorname{Ind}_\pi(C),\unsure{\mathcal E|E})$. Comme $g\mid C\simeq f$, on \deleted{\unsure{en}} en conclut que le foncteur (4.4.16.2) est ess. surjectif, ce qui achève de prouver 4.4.16. De plus, on a construit un foncteur quasi-inverse explicite, $\psi$,
% manuscript position: foot
% archive sheet number: 188
